\documentclass{situationist-preprint}

\title{Epicurus Answered}
\msid{EXI-02-b}
\msfamily{EXI-02}
\msdiscipline{Meaning, Death and the Absurd}
\mstype{Companion manuscript}
\mskeywords{Epicurus; death; deprivation account; fear of death; Lucretius}

\begin{document}
\maketitle

\begin{abstract}
Epicurus argued that death is nothing to us, and drew from this that it should not be feared. EXI-02-a answers the first claim: death is bad for the one who dies because it deprives her of goods, and deprivation does not require a subject at the time of loss. We reconstruct Epicurus's argument as a chain of two inferences and show where the answer bears on it. The inference to the conclusion that death is not bad fails; the inference from badness to fear is a separate matter, and its failure would not rescue Epicurus. Epicurus is right that death should not be feared for a reason he did not give, and wrong that it is nothing to us.
\end{abstract}

\section{The argument reconstructed}

In the \emph{Letter to Menoeceus}, Epicurus argues that death is nothing to us, since good and evil lie in sensation and death is the privation of sensation, and that when we are, death is not present, and when death is present, we are not \citep{EpicurusMen}. He concludes that the fear of death is groundless. The argument has two stages.

\begin{enumerate}
\item[(E1)] What is bad for a person must be experienced by her, or must at least occur while she exists. Death is neither. Hence death is not bad for the one who dies.
\item[(E2)] It is irrational to fear what is not bad for one. Hence it is irrational to fear death.
\end{enumerate}

\section{The answer}

EXI-02-a answers (E1). Its premise combines two requirements: an experience requirement and a requirement of co-existence with the loss. Lemma~3.1 of EXI-02-a shows the first to be false: some events are bad for a person though she never experiences them as bad \citep{Nagel1970}. Lemma~3.2 shows that the second, read as a demand that there be a subject of whom the comparative fact of the deprivation principle (Principle~P1 of EXI-02-a) holds, is satisfied, since the person who dies exists at the times she lives and stands tenselessly in the comparison (Assumption~2.3 of EXI-02-a). Read instead as a demand that the subject exist at the time at which the lost goods would have been, it is not a requirement on comparative badness at all. Theorem~3.3 and Corollary~3.4 of EXI-02-a conclude that death is bad for the one who dies in the ordinary case, and that (E1) fails for this reason.

\begin{proposition}\label{p:stage}
The failure of (E1) leaves (E2) to be assessed independently. If (E2) fails as well, it fails as an inference from badness to fear, not because death is after all nothing to us.
\end{proposition}

This matters because Epicurus's therapeutic aim, the removal of fear, can be pursued on other grounds. EXI-05 argues that fear of death is not rational, though regret at the prospect of death may be. If that is right, Epicurus's practical conclusion survives in part, and his metaphysical premise does not. Fear presents its object as a danger to be averted or escaped; death is a deprivation that, for the most part, cannot be averted and will not be suffered. The fitting response to an unavoidable deprivation is regret, not fear. Epicurus therefore reached a defensible attitude through an indefensible argument.

\begin{objection}
\citet{Rosenbaum1986} defends Epicurus by restricting the thesis to the state of being dead: being dead is not bad for anyone, whatever may be said of dying.
\end{objection}

\begin{reply}
The deprivation account agrees that being dead contains no evils for the one who is dead. Its claim is comparative and concerns the whole life: a life that ends earlier contains fewer goods than one that ends later \citep{Feldman1991}. The restricted thesis is therefore true and does not touch the badness of death; it only relocates the badness from a state to a comparison between lives.
\end{reply}

\section{Conclusion}

Epicurus's argument fails at its first step, because death deprives the one who dies and deprivation needs no subject at the time of loss. What survives of his teaching is the counsel against fear, which stands on separate grounds and leaves the badness of death in place.

\begin{thebibliography}{9}\small
\bibitem[Epicurus(Men.)]{EpicurusMen} Epicurus. \emph{Letter to Menoeceus}. In Diogenes Laertius, \emph{Lives of Eminent Philosophers} X.
\bibitem[Feldman(1991)]{Feldman1991} Feldman, F. (1991). Some puzzles about the evil of death. \emph{The Philosophical Review}, 100(2), 205--227.
\bibitem[Nagel(1970)]{Nagel1970} Nagel, T. (1970). Death. \emph{No\^us}, 4(1), 73--80.
\bibitem[Rosenbaum(1986)]{Rosenbaum1986} Rosenbaum, S.~E. (1986). How to be dead and not care: A defense of Epicurus. \emph{American Philosophical Quarterly}, 23(2), 217--225.
\end{thebibliography}

\end{document}
